\documentclass[10pt,a4paper]{amsart}
\usepackage[margin=19mm]{geometry}
\usepackage{amsmath,amssymb,mathtools}
\usepackage[scr=rsfs,cal=euler]{mathalfa}
\usepackage{xfrac}
\usepackage[unicode,hidelinks,bookmarks=false]{hyperref}
\newcommand{\ie}{i.e.\ }
\newcommand{\cf}{cf.\ }
\newcommand{\resp}{respectively\ }
\newcommand{\Subsection}{\subsection}
\providecommand{\nobreakdash}{\nobreak\hbox{-}}
\setlength{\emergencystretch}{2em}
\multlinegap0pt
\usepackage{bm}
\usepackage{bbm}
\usepackage{dsfont}
\usepackage{xspace}
\usepackage{cancel}
\usepackage{mathtools}
\usepackage{amsthm}
\makeatletter
\@ifundefined{thm}{\newtheorem{thm}{Thm}}{}
\@ifundefined{prop}{\newtheorem{prop}[thm]{Proposition}}{}
\makeatother
\makeatletter
\newcommand{\card}[1]{\left|\left\{ #1 \right\}\right|}
\@ifundefined{idproof}{\newenvironment{idproof}{
  \noindent \textit{Sketch of proof. }}{\hfill \(\square\) }}{}
\DeclareMathOperator{\inv}{inv}
\DeclareMathOperator{\noninv}{noninv}
\@ifundefined{preuve}{\newenvironment{preuve}{
  \noindent \textit{Proof. }}{\hfill \(\square\) }}{}
\makeatother
\title[Large deviation principles for pattern-avoiding permutations]{Large deviation principles for pattern-avoiding permutations, and limit shapes for constrained Mallows permutations}
\author{Thomas Budzinski, Victor Dubach, Valentin F\'eray, Mohamed Slim Kammoun, Myl\`ene Ma\"ida}
\date{Source: arXiv:2604.23751v1 (CC BY 4.0)}
\begin{document}
\maketitle

\noindent\emph{Source and licence.} Large deviation principles for pattern-avoiding permutations, and limit shapes for constrained Mallows permutations. Version arXiv:2604.23751v1 (CC BY 4.0), distributed under the Creative Commons Attribution 4.0 International licence, as recorded in the arXiv licence metadata for this exact version and shown on the abstract page https://arxiv.org/abs/2604.23751v1. Full rights notice: licenses/arxiv-2604.23751-CC-BY-4.0.txt.\par\smallskip

\noindent\textit{Editorial scope.} Counted unit: the Prop whose source label is \texttt{prop: inv of 231 avoiding permutation} in the article (source0.tex lines 1750--1757, source label \texttt{prop: inv of 231 avoiding permutation}), delivered together with the article's own complete original proof of it (source0.tex lines 1759--1814, one \texttt{proof} environment of 56 lines). No mathematics is written, repaired or re-derived: statement and proof are the article's own text line for line. Required context retained uncounted: none. Every definition, display and label of those lines is retained unchanged. Omitted from this package and not redistributed: 1--1749, 1758--1758, 1815--2593. Nothing inside a retained excerpt is omitted or altered: every retained body is delivered exactly as the article's lines are written, so no omission receipt of any kind accompanies this package. Citations: the retained excerpts carry no citation command at all, so no bibliography entry of the article is transcribed and no reference is redistributed. Reference adaptation: none was needed and none was made; every reference inside the delivered unit resolves inside it, so no unresolved reference or citation is produced. Printed slips kept literal: no printed word, symbol, hypothesis or number of the counted statement or of its proof was altered. Layout and class: the article's own class is \texttt{amsart} with packages including graphicx, amssymb, babel, inputenc, amsmath, amsfonts, amsthm, fontenc, fullpage, hyperref, enumitem, mathrsfs, relsize, tikz; the wrapper is the lane's pinned \texttt{amsart} editorial wrapper at 10pt on A4, which restates the theorem-like environments the retained text needs and adds the article's own macro definitions that the retained text needs (\texttt{\textbackslash card}, \texttt{\textbackslash inv}, \texttt{\textbackslash noninv}), with extra wrapper packages (bm, bbm, dsfont, xspace, cancel, mathtools). The delivered numbering is the unit's own. Assets: no retained span contains a figure environment, a graphics-inclusion command, a PDF-page inclusion, a table or a TikZ picture; no image file is copied into this package, the package declares no asset dependency and no image file is redistributed with it. Licence: version arXiv:2604.23751v1 of this article is distributed under the Creative Commons Attribution 4.0 International licence, as recorded in the arXiv licence metadata for this exact version and shown on the abstract page https://arxiv.org/abs/2604.23751v1. A full rights notice is delivered at licenses/arxiv-2604.23751-CC-BY-4.0.txt. Characters: every retained excerpt is pure ASCII, so the delivered engine, XeLaTeX with its default Latin Modern text font, reports no missing character.
\begin{prop}\label{prop: inv of 231 avoiding permutation}
    If $\sigma$ is a $231$-avoiding permutation then
    \[\inv(\sigma) = \binom{n}{2} - \sum_{x=0}^{n-1} F_\sigma(x).\]
    Therefore,
    \[ \inv(\sigma)
    = \binom{n}{2} - n^2 \int_0^1 f_\sigma(x) \mathrm{d}x 
    = \frac{-n}{2} + 2 n^2 \int_0^1 \varphi_\sigma(x) \mathrm{d}x  .\]
\end{prop}

\begin{proof}
 Let $\noninv(\sigma) := \card{ (x,y)\in[n]^2 : x<y , \sigma(x)<\sigma(y) }$ be the number of noninversions of $\sigma$.
    Since $\inv(\sigma) + \noninv(\sigma) = \binom{n}{2}$, we may as well compute $\noninv(\sigma)$.
    For this, we enumerate noninversions by their north-east point.
    Let $c\in[n]$ be fixed.
    We claim that 
    \begin{equation}\label{eq: noninversion set of 231 avoiding}
        \left\{ x\in[n] : x<c , \sigma(x)<\sigma(c) \right\}
        = \left\{ \sigma^{-1}(y) : y \le F_\sigma(c-1) \right\} \,,
    \end{equation}
    and thus
    \begin{equation*}\label{eq: noninversion count of 231 avoiding}
        \card{ x\in[n] : x<c , \sigma(x)<\sigma(c) }
        = F_\sigma(c-1) \,,
    \end{equation*}
    so that $\noninv(\sigma) =  \sum_{x=0}^{n-1} F_\sigma(x).$

Indeed, for $c\in [n],$ by definition of $F_\sigma,$ for any $y \le    F_\sigma(c-1)$, we have $\sigma^{-1}(y)<c $ since there is no point under the RLM curve.

Moreover, by construction $\sigma(c) \ge  F_\sigma(c-1) +1.$ If $ \sigma(c) =  F_\sigma(c-1) +1,$ then the previous remark directly gives \eqref{eq: noninversion set of 231 avoiding}. Now, if  $ \sigma(c) > F_\sigma(c-1) +1,$ then there exists $c'>c$ such that  $ \sigma(c') = F_\sigma(c-1) +1$ (intuitively, $c'$ is the first RL minimum at the right of $c$). Let $ F_\sigma(c-1)< y< \sigma(c).$
If $\sigma^{-1}(y) < c,$ then the permutation $\sigma$ is not $231$-avoiding because we have $\sigma^{-1}(y) < c < c'$ and
$\sigma(c')< y < \sigma(c).$ Therefore, \eqref{eq: noninversion set of 231 avoiding} holds in both cases and this concludes the proof of the first formula.

%     Indeed, if $c = a_{i(c)}$, then \eqref{eq: noninversion set of 231 avoiding} follows directly from the fact that $(c, \sigma(c)) = \left(a_{i(c)}, b_{i(c)}\right)$ is an RL minimum.
%     Otherwise, assume that $c\notin\left\{ a_i : i\in[\ell] \right\}$, \textit{i.e.}~that $c < a_{i(c)}$.
%
%     On the one hand, let $x\in[n]$ be such that $\sigma(x) < b_{i(c)}$.
%     Since $\left( a_{i(c)}, b_{i(c)} \right)$ is an RL minimum, $x<a_{i(c)}$.
%     Furthermore $x\notin \left( a_{i(c)-1}, a_{i(c)} \right)$, because otherwise $(x,\sigma(x))$ would be an RL minimum.
%     Therefore, $x<c$.
%     Since $c\in\left( a_{i(c)-1}, a_{i(c)} \right)$, it also holds that $\sigma(c) > b_{i(c)}$, and thus $\sigma(x) < \sigma(c)$.
%     This shows that the RHS of \eqref{eq: noninversion set of 231 avoiding} is contained in the LHS.
%
%     On the other hand, let $x<c$ be such that $\sigma(x)<\sigma(c)$.
%     As before, we have that $\sigma(c) > b_{i(c)}$.
%     Then $x< c< a_{i(c)}$ are such that $\sigma(c) > \sigma(x)$ and $\sigma(c) > b_{i(c)}$.
%     Since $\sigma$ avoids $231$, we get that $\sigma(x) < b_{i(c)}$.
%     This proves \eqref{eq: noninversion set of 231 avoiding}, and thus \eqref{eq: noninversion count of 231 avoiding}.
%     Hence:
%     \begin{equation*}
%         \noninv(\sigma)
%         = \sum_{c=1}^n b_{i(c)}-1
%         = \sum_{i=1}^\ell \left( a_i-a_{i-1} \right) \left( b_i-1 \right)
%     \end{equation*}
%     with the convention $a_0=0$.
%     Since $a_\ell=n$, we deduce that:
%     \begin{equation*}
%         \noninv(\sigma)
%         = \left( \sum_{i=1}^\ell \left( a_i-a_{i-1} \right) b_i \right) - n \,.
%     \end{equation*}
%     Since $F_\sigma(c-1) = b_{i(c)}-1$ for each $c\in[n]$, we can also write $\noninv(\sigma) = \sum_{x=0}^{n-1} F_\sigma(x)$.

    Finally, since the non-normalized RLM curve $F_\sigma$ is locally constant and jumps at integer positions,  
    $\sum_{x=0}^{n-1} F_\sigma(x)=\int_0^n F_\sigma(x) \mathrm{d}x$.
    Using $\int_0^1 f_\sigma(x) \mathrm{d}x = \frac{1}{n^2} \int_0^n F_\sigma(x) \mathrm{d}x$ and $\frac12 = \int_0^1 f_\sigma(x) \mathrm{d}x + 2\int_0^1 \varphi_\sigma(x) \mathrm{d}x$, the rest of the Proposition follows.
\end{proof}

\end{document}
